\begin{proofbox}
  \textbf{\textcolor{CProof}{证明一：看图像（最好理解）}}

  \textbf{想法：} 恒正就是图像全部在 $x$ 轴上方。
  \begin{enumerate}[leftmargin=*,itemsep=0.3em]
    \item 若开口向下（$a<0$），两端会越走越低，不可能恒正；所以必须 $a>0$。
    \item 开口向上以后，还得\textbf{完全不碰 $x$ 轴}。$f(x)=0$ 无实数解，等价于 $\Delta<0$。
    \item 反过来，$a>0$ 且 $\Delta<0$，就是开口向上、没有任何横轴交点，图像必在横轴上方，因此恒正。
  \end{enumerate}
  \begin{center}
  \begin{tikzpicture}[x=0.95cm,y=0.6cm,>=stealth]
    \draw[->,gray] (-1.9,0)--(1.95,0) node[right] {$x$};
    \draw[->,gray] (0,-3.25)--(0,3.3) node[above] {$y$};
    \draw[green!55!black,very thick,domain=-1.5:1.5,samples=50] plot (\x,{\x*\x+0.5});
    \draw[red!75!black,very thick,domain=-1.5:1.5,samples=50] plot (\x,{-\x*\x-0.5});
    \node[anchor=west,green!45!black] at (2.15,2) {$a>0,\ \Delta<0$：恒正};
    \node[anchor=west,red!65!black] at (2.15,-1.5) {$a<0,\ \Delta<0$：恒负};
  \end{tikzpicture}

  \end{center}
  \textbf{提醒：} 只有 $\Delta<0$ 不够，红色曲线也没有交点，却全部在轴下方。

  \medskip
  \textbf{\textcolor{CProof}{证明二：配方（每一步都能算出来）}}

  \textbf{想法：} 把式子变成“一个平方 $+$ 一个固定数”，看最低能到多少。
  \[
    \begin{aligned}
      ax^2+bx+c
      &=a\left(x^2+\frac ba x\right)+c\\
      &=a\left(x+\frac{b}{2a}\right)^2+c-\frac{b^2}{4a}\\
      &=a\left(x+\frac{b}{2a}\right)^2-\frac{\Delta}{4a}.
    \end{aligned}
  \]
  \begin{enumerate}[leftmargin=*,itemsep=0.3em]
    \item 恒正必有 $a>0$：否则 $a<0$，平方项取很大时，整个式子就会变成负数。
    \item 在 $a>0$ 下，平方项最小为 $0$，发生在 $x=-\dfrac b{2a}$；因此最低值为 $-\dfrac\Delta{4a}$。
    \item 恒正 $\Longleftrightarrow$ 最低值 $>0$ $\Longleftrightarrow -\dfrac\Delta{4a}>0$ $\Longleftrightarrow \Delta<0$。
  \end{enumerate}
  \begin{center}
  \begin{tikzpicture}[x=1cm,y=0.7cm,>=stealth]
    \draw[->,gray] (-1.05,0)--(3.1,0) node[right] {$x$};
    \draw[->,gray] (0,-0.4)--(0,3.45) node[above] {$y$};
    \draw[green!55!black,very thick,domain=-0.45:2.45,samples=50] plot (\x,{(\x-1)*(\x-1)+1});
    \draw[dashed,gray] (1,0)--(1,1);
    \fill[green!55!black] (1,1) circle (2pt);
    \node[anchor=west] at (3.75,.4) {最低点 $(1,1)$};
    \node[anchor=west] at (3.75,2.6) {$f(x)=(x-1)^2+1$};
    \node[anchor=west] at (3.75,1.5) {平方部分 $\ge0$};
  \end{tikzpicture}

  \end{center}
  \textbf{小例子：} $(x-1)^2+1\ge1>0$，你已经亲手用配方证明了一个恒正式子。

  \medskip
  \textbf{\textcolor{CProof}{证明三：顶点最小值（考场最快）}}

  \textbf{想法：} 开口向上的抛物线，最低处就是顶点。
  对 $a>0$，顶点横坐标是 $x=-\dfrac{b}{2a}$，代入直接得
  \[
     f_{\min}=f\left(-\frac{b}{2a}\right)=\frac{4ac-b^2}{4a}=-\frac{\Delta}{4a}.
  \]
  因此
  \[
    f(x)>0\ \text{恒成立}\quad\Longleftrightarrow\quad
    \underbrace{a>0}_{\text{开口向上}}\quad\text{且}\quad
    \underbrace{f_{\min}>0}_{\text{最低点在轴上方}}\quad\Longleftrightarrow\quad
    a>0,\ \Delta<0.
  \]
  \begin{center}
  \begin{tikzpicture}[x=1.1cm,y=0.63cm,>=stealth]
    \draw[->,gray] (-1.8,0)--(1.85,0) node[right] {$x$};
    \draw[->,gray] (0,-1.45)--(0,3) node[above] {$y$};
    \draw[green!55!black,very thick,domain=-1.35:1.35,samples=45] plot (\x,{\x*\x+0.75});
    \draw[orange!85!black,thick,domain=-1.62:1.62,samples=45] plot (\x,{\x*\x});
    \draw[red!75!black,thick,domain=-1.65:1.65,samples=45] plot (\x,{\x*\x-0.75});
    \fill[green!55!black] (0,.75) circle (1.7pt);
    \fill[orange!85!black] (0,0) circle (1.7pt);
    \fill[red!75!black] (0,-.75) circle (1.7pt);
    \node[anchor=west,green!45!black] at (2.05,1.5) {最低值 $>0$：恒正};
    \node[anchor=west,orange!80!black] at (2.05,.7) {最低值 $=0$：仅非负};
    \node[anchor=west,red!70!black] at (2.05,-.1) {最低值 $<0$：不恒正};
  \end{tikzpicture}

  \end{center}
  \textbf{何时选它：} 已知顶点式或很容易配方，就直接看最低值；一般含参数的标准二次式，直接检验 $a>0$ 且 $\Delta<0$ 更快。
\end{proofbox}
